\chapter{CƠ SỞ LÝ THUYẾT VÀ CÁC NGHIÊN CỨU LIÊN QUAN}
\label{chap:co-so-ly-thuyet}

\section{Phân bổ tài nguyên và nhiễu trong mạng vô tuyến}
\label{sec:phan-bo-tai-nguyen}

Phân bổ tài nguyên trong mạng vô tuyến là quá trình quyết định cách sử dụng các tài
nguyên hữu hạn như công suất phát, thời gian, tần số, mã, không gian và các tham số điều
khiển của môi trường truyền. Trong hệ thống nhiều người dùng, một quyết định tốt không
chỉ làm tăng tốc độ của một liên kết riêng lẻ mà còn phải xem xét tác động của quyết
định đó lên các liên kết khác. Đây là nguồn gốc của bài toán quản lý nhiễu.

Trong một hệ thống có nhiều UE, cùng một đơn vị công suất có thể mang lại lợi ích rất
khác nhau tùy theo kênh và mức nhiễu hiện tại. Nếu dồn phần lớn công suất cho UE có kênh
tốt, tổng tốc độ có thể tăng nhanh nhưng UE yếu dễ vi phạm yêu cầu tối thiểu. Nếu chia
công suất đều để tăng tính công bằng, nhiều luồng cùng tồn tại có thể tạo nhiễu lẫn nhau
và làm giảm hiệu suất sử dụng công suất. Vì vậy, bài toán thực tế thường là một bài toán
cân bằng giữa ba yêu cầu: hiệu năng tổng thể, QoS của từng UE và chi phí tính toán của
bộ điều khiển.

Trong luận văn này, chỉ số hiệu năng chính là tổng tốc độ. Tuy nhiên, tổng tốc độ không
được dùng một cách tách rời. Một nghiệm có tổng tốc độ cao nhưng để một hoặc nhiều UE có
tốc độ gần bằng không không được xem là nghiệm phù hợp với mục tiêu phục vụ đa người
dùng. Do đó, QoS được đưa vào cả phát biểu bài toán và hàm phần thưởng dùng cho DRL.
Cách tiếp cận này phản ánh một yêu cầu vật lý đơn giản: hệ thống chỉ hữu ích khi không
chỉ truyền nhanh mà còn duy trì mức dịch vụ tối thiểu cho các UE đang được phục vụ.

\section{Các kỹ thuật đa truy cập}
\label{sec:ky-thuat-da-truy-cap}

\subsection{Từ OMA đến SDMA và NOMA}
\label{ssec:oma-sdma-noma}

Đa truy cập là cơ chế cho phép nhiều UE cùng chia sẻ tài nguyên của một BS. OMA tách
người dùng trên các miền trực giao như thời gian hoặc tần số. Ưu điểm của cách này là
nhiễu liên người dùng trong cùng tài nguyên được giảm mạnh và bộ thu tương đối đơn giản.
Nhược điểm là một phần tài nguyên phải dành riêng cho từng UE, nên hiệu suất phổ có thể
giảm khi số UE tăng hoặc khi yêu cầu tốc độ biến động nhanh.

SDMA tận dụng miền không gian của hệ thống nhiều anten. BS sử dụng các véc-tơ tiền mã
hóa để tập trung năng lượng theo các hướng khác nhau và giảm nhiễu chéo. Khi số anten đủ
lớn và kênh giữa các UE đủ phân biệt, SDMA có thể phục vụ đồng thời nhiều luồng trên
cùng thời gian--tần số. Tuy nhiên, lợi thế này phụ thuộc trực tiếp vào khả năng tạo phân
tách không gian. Trong hệ SISO của luận văn, BS chỉ có một anten nên không tồn tại véc-tơ
beamforming nhiều chiều để tạo các hướng truyền riêng cho từng UE. Đây là khác biệt quan
trọng khi phân tích nhiễu của các luồng riêng ở Chương~\ref{chap:mo-hinh-he-thong}.

NOMA cho phép nhiều UE sử dụng cùng tài nguyên thời gian--tần số nhưng phân biệt bằng
miền công suất. Một UE có thể giải mã tín hiệu của UE khác trước rồi loại bỏ bằng SIC.
Cách tiếp cận này hữu ích khi kênh của các UE chênh lệch rõ, nhưng thứ tự giải mã, phân
bổ công suất và lan truyền sai số SIC có thể làm bài toán thiết kế phức tạp. Các nghiên
cứu tổng quan cũng chỉ ra rằng không nên xem NOMA như một cơ chế mặc nhiên vượt trội
trong mọi cấu hình nhiều anten hoặc mọi phân bố kênh \cite{Clerckx2021NOMACritical}.

Hình~\ref{fig:ch1-da-truy-cap} minh họa trực giác của bốn cách tổ chức tín hiệu. Hình
không nhằm mô tả đầy đủ mọi biến thể chuẩn hóa, mà tập trung vào cách mỗi kỹ thuật xử lý
việc nhiều người dùng cùng chia sẻ tài nguyên.

\begin{figure}[htbp]
\centering
\resizebox{\textwidth}{!}{\begin{tikzpicture}[font=\scriptsize,>=Stealth,line width=0.5pt]

% =====================================================================
% (a) OMA
% =====================================================================
\begin{scope}[shift={(0,0)}]
  \node[tit] at (0,3.95) {(a) OMA — tách tài nguyên};
  \foreach \x/\y/\lbl/\c in {0.5/0.7/UE$_3$/blue, 2.2/0.7/UE$_4$/violet,
                             0.5/1.7/UE$_1$/orange, 2.2/1.7/UE$_2$/teal}{
     \fill[\c!18] (\x,\y) rectangle ++(1.7,1.0);
     \draw[\c!60!black,line width=0.6pt] (\x,\y) rectangle ++(1.7,1.0);
     \node[font=\scriptsize] at (\x+0.85,\y+0.5) {\lbl};
  }
  \draw[gray!60,->] (0.5,0.5) -- (4.1,0.5) node[anchor=north,font=\tiny,gray!50!black,midway,yshift=-2pt]{thời gian};
  \draw[gray!60,->] (0.35,0.7) -- (0.35,2.9) node[anchor=east,font=\tiny,gray!50!black,midway,rotate=90,yshift=6pt]{tần số};
  \node[note] at (0,0.15)
    {Mỗi UE chiếm một khối tài nguyên riêng. Không có nhiễu liên người dùng,
     nhưng mỗi UE chỉ dùng được $1/K$ tài nguyên.};
\end{scope}

% =====================================================================
% (b) SDMA
% =====================================================================
\begin{scope}[shift={(8.6,0)}]
  \node[tit] at (0,3.95) {(b) SDMA — tách theo không gian};
  % BS nhieu anten
  \node[bx,minimum height=0.9cm,minimum width=0.5cm] (bs) at (0.75,1.7) {BS};
  \foreach \dy in {-0.25,0,0.25}{
     \draw[blue!70!black,line width=0.6pt] (0.98,1.7+\dy) -- ++(0.22,0);
     \fill[blue!65] (1.20,1.7+\dy-0.07) -- (1.20,1.7+\dy+0.07) -- (1.36,1.7+\dy) -- cycle;
  }
  % 4 bup song
  \foreach \y/\lbl/\c in {2.90/UE$_1$/orange, 2.20/UE$_2$/teal, 1.20/UE$_3$/blue, 0.50/UE$_4$/violet}{
     \fill[\c!22,opacity=0.75] (1.40,1.7) -- (4.15,\y+0.30) -- (4.15,\y-0.30) -- cycle;
     \node[ue,fill=\c!20,draw=\c!60!black] at (4.45,\y) {\lbl};
  }
  \node[note] at (0,0.15)
    {BS dùng véc-tơ tiền mã hóa để tạo búp sóng riêng cho từng UE.
     \textbf{Cần nhiều anten.} Hệ SISO của luận văn chỉ có một anten nên
     không tạo được phân tách không gian.};
\end{scope}

% =====================================================================
% (c) NOMA
% =====================================================================
\begin{scope}[shift={(0,-5.6)}]
  \node[tit] at (0,3.95) {(c) NOMA — tách theo miền công suất};
  % chong cong suat
  \foreach \y/\h/\lbl/\c in {0.75/0.95/$p_1$/orange, 1.70/0.70/$p_2$/teal,
                             2.40/0.48/$p_3$/blue, 2.88/0.32/$p_4$/violet}{
     \fill[\c!22] (0.45,\y) rectangle ++(0.75,\h);
     \draw[\c!60!black,line width=0.5pt] (0.45,\y) rectangle ++(0.75,\h);
     \node[font=\tiny] at (0.825,\y+\h/2) {\lbl};
  }
  \node[font=\tiny,anchor=north,align=center] at (0.825,0.70) {chồng\\ công suất};
  \draw[gray!60,->] (0.30,0.75) -- (0.30,3.25) node[anchor=east,font=\tiny,gray!50!black,midway,rotate=90,yshift=6pt]{công suất};

  % chuoi SIC nhieu tang: mot cot doc, khong cat nhau
  \node[sm,anchor=west,minimum width=1.6cm] (n1) at (2.05,3.05) {Giải mã $s_1$};
  \node[sm,anchor=west,minimum width=1.6cm,fill=gray!22] (n2) at (2.05,2.55) {SIC};
  \node[sm,anchor=west,minimum width=1.6cm] (n3) at (2.05,2.05) {Giải mã $s_2$};
  \node[sm,anchor=west,minimum width=1.6cm,fill=gray!22] (n4) at (2.05,1.55) {SIC};
  \node[sm,anchor=west,minimum width=1.6cm] (n5) at (2.05,1.05) {Giải mã $s_k$};
  \foreach \a/\b in {n1/n2, n2/n3, n3/n4, n4/n5}{ \draw[->] (\a.south) -- (\b.north); }
  \draw[decorate,decoration={brace,amplitude=4pt},gray!70] (3.80,3.20) -- (3.80,0.92);
  \node[font=\tiny,anchor=west,align=left,gray!40!black] at (3.98,2.06) {tối đa\\ $K-1$\\ tầng SIC};

  \node[note] at (0,0.18)
    {Các UE dùng chung tài nguyên, phân biệt bằng mức công suất.
     UE kênh tốt phải giải mã và loại lần lượt tín hiệu của các UE khác.
     Thứ tự giải mã phụ thuộc độ mạnh kênh.};
\end{scope}

% =====================================================================
% (d) RSMA
% =====================================================================
\begin{scope}[shift={(8.6,-5.6)}]
  \node[tit] at (0,3.95) {(d) RSMA — chung + riêng};
  % tach thong diep
  \node[font=\tiny,anchor=east] at (0.42,2.95) {$W_k$};
  \foreach \i/\x/\c in {1/0.50/orange, 2/1.42/teal, 3/2.34/blue, 4/3.26/violet}{
     \fill[magenta!25] (\x,2.80) rectangle ++(0.34,0.32);
     \draw[magenta!70!black,line width=0.4pt] (\x,2.80) rectangle ++(0.34,0.32);
     \fill[\c!25] (\x+0.34,2.80) rectangle ++(0.48,0.32);
     \draw[\c!60!black,line width=0.4pt] (\x+0.34,2.80) rectangle ++(0.48,0.32);
     \node[font=\tiny,anchor=north] at (\x+0.41,2.74) {\i};
  }
  \node[font=\tiny,anchor=west,magenta!70!black] at (4.20,3.08) {phần chung $W_{c,k}$};
  \node[font=\tiny,anchor=west,gray!40!black] at (4.20,2.80) {phần riêng $W_{p,k}$};

  % gop thanh luong chung + cac luong rieng
  \node[sm,fill=magenta!15,anchor=west] (sc) at (0.50,1.95) {gộp $\to$ luồng chung $s_c$};
  \node[sm,anchor=west] (sp) at (3.45,1.95) {$s_1,\dots,s_4$};
  \draw[->] (1.00,2.72) -- (1.00,2.22);
  \draw[->] (4.00,2.72) -- (4.00,2.22);

  % may thu cua UE k : dung 1 tang SIC
  \node[sm,anchor=west] (r1) at (0.50,1.05) {Giải mã $s_c$};
  \node[sm,fill=magenta!12,anchor=west] (r2) at (2.15,1.05) {SIC (1 tầng)};
  \node[sm,anchor=west] (r3) at (4.00,1.05) {Giải mã $s_k$};
  \draw[->] (r1.east) -- (r2.west);
  \draw[->] (r2.east) -- (r3.west);
  \draw[->] (1.20,1.75) -- (1.20,1.28);
  \draw[->] (4.30,1.75) -- (4.30,1.28);
  \node[font=\tiny,anchor=north west,align=left] at (0.50,0.90)
       {$R_k=\underbrace{\eta_k R_c}_{\text{phần chung}}+\underbrace{R_{p,k}}_{\text{phần riêng}}$};
  \node[note] at (0,0.15)
    {Thông điệp mỗi UE được tách thành phần chung và phần riêng.
     \textbf{Mọi UE giải luồng chung trước, rồi dùng đúng một tầng SIC}
     để loại $s_c$, sau đó giải luồng riêng của mình.};
\end{scope}

% =====================================================================
% Dai quan he bao ham
% =====================================================================
\begin{scope}[shift={(0,-8.05)}]
  \draw[gray!45,line width=0.9pt,->] (1.6,0) -- (12.2,0);
  \foreach \x in {2.0,11.8}{ \draw[gray!60] (\x,-0.12) -- (\x,0.12); }
  \node[font=\tiny,anchor=north] at (2.0,-0.18) {$\rho_c\to 0$};
  \node[font=\tiny,anchor=north] at (11.8,-0.18) {$\rho_c\to 1$};
  \node[bx,fill=teal!12,anchor=north] at (2.0,-0.55) {SDMA\\ \tiny chỉ luồng riêng};
  \node[bx,fill=violet!12,anchor=north] at (11.8,-0.55) {NOMA\\ \tiny gần như chỉ luồng chung};
  \fill[magenta!14] (5.0,-0.30) rectangle (9.0,0.30);
  \draw[magenta!60!black,line width=0.6pt] (5.0,-0.30) rectangle (9.0,0.30);
  \node[font=\scriptsize,magenta!60!black] at (7.0,0) {RSMA};
  \node[font=\tiny,anchor=north,align=center] at (7.0,-0.55)
       {RSMA bao hàm SDMA và NOMA như hai trường hợp đặc biệt;\\
        luận văn giới hạn miền hành động ở $\rho_c\in[0{,}70;\,0{,}99]$};
  \node[font=\tiny,anchor=east] at (1.5,0) {$\rho_c=p_c/P_{\max}$:};
\end{scope}

\end{tikzpicture}
}
\caption{So sánh trực giác của OMA, SDMA, NOMA và RSMA với $K=4$ người dùng}
\label{fig:ch1-da-truy-cap}
\figuresource{Tác giả xây dựng dựa trên \cite{Mao2018RSMABridging, Mao2022RSMAFundamentals, Clerckx2021NOMACritical} (2026)}
\end{figure}

\subsection{Nguyên lý của RSMA}
\label{ssec:nguyen-ly-rsma}

RSMA tiếp cận quản lý nhiễu theo cách mềm hơn. Thông điệp dành cho mỗi UE được chia thành
hai phần: một phần được gộp vào luồng chung và một phần nằm trong luồng riêng. Mọi UE đều
giải mã luồng chung trước trong khi xem tất cả luồng riêng như nhiễu. Sau khi luồng chung
được giải mã thành công, UE dùng SIC để loại luồng chung khỏi tín hiệu thu. Tiếp theo, UE
giải mã luồng riêng của mình và xem các luồng riêng của UE khác như nhiễu
\cite{Mao2018RSMABridging, Mao2022RSMAFundamentals}.

Cách tổ chức này tạo ra một công cụ điều chỉnh linh hoạt. Nếu một phần giao thoa khó xử
lý khi giữ hoàn toàn ở luồng riêng, bộ điều khiển có thể chuyển một phần thông tin sang
luồng chung để nhiều UE cùng giải mã. Ngược lại, nếu kênh của một UE thuận lợi, phần
riêng của UE đó có thể được tăng để khai thác tốc độ cao. Trong luận văn, biến phân bổ
tốc độ chung đóng vai trò quan trọng vì nó quyết định UE nào nhận bao nhiêu phần của tốc
độ luồng chung sau khi tốc độ chung đã được xác định bởi UE khó giải mã nhất.

Một điểm cần nhấn mạnh là luồng chung không phải một luồng dữ liệu ``phụ'' có thể bỏ qua.
Trong bài toán QoS, luồng chung có thể đóng vai trò như phần tốc độ bảo đảm cho các UE
yếu. Nếu một UE không được cấp nhiều công suất riêng, UE đó vẫn có thể đạt tốc độ tối
thiểu nhờ phần tốc độ chung được phân bổ cho nó. Đây là cơ sở trực giác cho việc tỷ lệ
công suất luồng chung cao có thể hữu ích trong cấu hình SISO đang xét.

\section{RIS và STAR--RIS}
\label{sec:ris-star-ris}

\subsection{RIS phản xạ truyền thống}
\label{ssec:ris-phan-xa}

RIS là một bề mặt gồm nhiều phần tử có thể điều chỉnh đáp ứng điện từ. Trong mô hình lý
tưởng thường dùng ở mức hệ thống, mỗi phần tử chủ yếu thay đổi pha của tín hiệu phản xạ.
Khi các pha được phối hợp hợp lý, các thành phần tín hiệu đi qua nhiều phần tử có thể
cộng tăng cường tại vị trí mong muốn. Nhờ đó, RIS thay đổi kênh hiệu dụng mà không cần
một bộ phát tích cực tại bề mặt \cite{Wu2019IRSBeamforming, DiRenzo2020SmartRadio,
Wu2021IRSTutorial}.

Tác dụng vật lý của RIS có thể hình dung như việc điều khiển nhiều đường truyền nhỏ sao
cho các ``mũi tên pha'' quay về gần cùng hướng tại UE. Nếu các thành phần đến ngược pha,
chúng triệt tiêu nhau và độ lợi kênh giảm. Nếu được căn pha, chúng cộng xây dựng và độ
lợi kênh tăng. Khái niệm căn pha này được sử dụng lại trong
Chương~\ref{chap:phuong-phap} để tạo pha tham chiếu cho bộ giải mã hành động.

\subsection{STAR--RIS và phạm vi phủ sóng hai phía}
\label{ssec:star-ris}

RIS phản xạ chỉ điều khiển đáng kể nửa không gian phía phản xạ. STAR--RIS được nghiên
cứu nhằm mở rộng vùng phủ bằng cách cho phép một phần năng lượng truyền qua bề mặt và
phần còn lại phản xạ \cite{Xu2021STARRIS, Mu2021STARAided, Khalid2022STARDesign}. Một
phần tử STAR--RIS trong mô hình ES có hai hệ số phức: một cho phía truyền và một cho phía
phản xạ.
\begin{equation}
\label{eq:phi-star}
\phi^T_n = \sqrt{\beta^T_n}\, e^{j\theta^T_n},
\qquad
\phi^R_n = \sqrt{\beta^R_n}\, e^{j\theta^R_n}.
\end{equation}

Trong đó: $\phi^T_n$ và $\phi^R_n$ là hệ số phức của phần tử thứ $n$ ở phía truyền và
phía phản xạ; $\beta^T_n$, $\beta^R_n$ là tỷ lệ năng lượng; $\theta^T_n$, $\theta^R_n$ là
pha; $j$ là đơn vị ảo.

Công thức~\eqref{eq:phi-star} tách tác động của một phần tử thành hai phần dễ hiểu: căn
bậc hai của $\beta$ quyết định biên độ tín hiệu đi về mỗi phía, còn $\theta$ quyết định
hướng pha của tín hiệu. Nhờ thay đổi hai nhóm biến này, STAR--RIS có thể vừa ``chia'' năng
lượng giữa hai phía vừa ``xoay'' pha để các đường truyền cộng thuận lợi tại UE.

Trong mô hình ES lý tưởng, năng lượng được bảo toàn theo quan hệ
$\beta^T_n + \beta^R_n = 1$. Luận văn dùng một biến $\beta^T_n$ và suy ra
$\beta^R_n = 1 - \beta^T_n$. Cách này vừa giảm số biến cần học vừa bảo đảm ràng buộc vật
lý ngay tại bộ giải mã hành động. Ngoài ES, tài liệu STAR--RIS còn mô tả các cơ chế
chuyển chế độ và chuyển theo thời gian; tuy nhiên các cơ chế đó không nằm trong phạm vi
mô phỏng của luận văn \cite{Mu2021STARAided, Liu2021STARCoupledPhase,
Wang2022STAROptimization}.

\begin{figure}[htbp]
\centering
\resizebox{\textwidth}{!}{\begin{tikzpicture}[font=\scriptsize,>=Stealth,line width=0.5pt]

% =====================================================================
% (a) RIS phan xa truyen thong
% =====================================================================
\begin{scope}[shift={(0,0)}]
  \node[tit] at (0,4.55) {(a) RIS phản xạ truyền thống};

  % vung khuat phia sau be mat
  \fill[gray!10] (3.30,0.30) rectangle (6.30,3.55);
  \draw[gray!45,dashed,line width=0.5pt] (3.30,0.30) rectangle (6.30,3.55);
  \node[font=\tiny,gray!45!black,align=center,anchor=north] at (4.80,1.95)
       {Nửa không gian phía sau\\ không được phục vụ};

  \risbar{3.18}{blue}
  \node[font=\tiny,anchor=north] at (3.18,0.22) {RIS, $N$ phần tử};

  % BS SISO
  \node[bx,minimum height=0.8cm,minimum width=0.5cm] (bs) at (0.45,1.85) {BS};
  \draw[blue!70!black,line width=0.6pt] (0.68,1.85) -- ++(0.20,0);
  \fill[blue!65] (0.88,1.78) -- (0.88,1.92) -- (1.04,1.85) -- cycle;

  % UE phia phan xa
  \node[ue,fill=orange!22,draw=orange!70!black] (u1) at (1.55,3.20) {UE$_1$};
  \node[ue,fill=blue!22,draw=blue!70!black]     (u3) at (1.75,0.62) {UE$_3$};
  % UE phia sau -- khong duoc phuc vu
  \node[ue,fill=gray!12,draw=gray!55,dashed] (u2) at (5.55,3.10) {UE$_2$};
  \node[ue,fill=gray!12,draw=gray!55,dashed] (u4) at (5.75,0.72) {UE$_4$};

  % tia toi va tia phan xa
  \draw[orange!80!black,->] (1.06,1.85) -- (3.05,2.72);
  \draw[orange!80!black,->] (3.05,2.72) -- (1.80,3.10);
  \draw[blue!70!black,->]   (1.06,1.85) -- (3.05,1.05);
  \draw[blue!70!black,->]   (3.05,1.05) -- (1.98,0.75);

  \node[font=\tiny,anchor=east,gray!45!black] at (3.05,3.74) {phía phản xạ};
  \node[font=\tiny,anchor=west,gray!45!black] at (3.32,3.74) {phía sau bề mặt};

  \node[font=\scriptsize,anchor=north west] at (0,-0.30)
       {Hệ số phần tử: \ $\phi_n=e^{j\theta_n}$};
\end{scope}

% =====================================================================
% (b) STAR-RIS
% =====================================================================
\begin{scope}[shift={(8.2,0)}]
  \node[tit] at (0,4.55) {(b) STAR--RIS};

  \fill[green!8] (3.30,0.30) rectangle (6.30,3.55);
  \draw[green!45!black,dashed,line width=0.5pt] (3.30,0.30) rectangle (6.30,3.55);

  \risbar{3.18}{green!60!black}
  \node[font=\tiny,anchor=north] at (3.18,0.22) {STAR--RIS, $N$ phần tử};

  \node[bx,minimum height=0.8cm,minimum width=0.5cm] (bs2) at (0.45,1.85) {BS};
  \draw[blue!70!black,line width=0.6pt] (0.68,1.85) -- ++(0.20,0);
  \fill[blue!65] (0.88,1.78) -- (0.88,1.92) -- (1.04,1.85) -- cycle;

  \node[ue,fill=orange!22,draw=orange!70!black] (v1) at (1.55,3.20) {UE$_1$};
  \node[ue,fill=blue!22,draw=blue!70!black]     (v3) at (1.75,0.62) {UE$_3$};
  \node[ue,fill=teal!22,draw=teal!70!black]     (v2) at (5.55,3.10) {UE$_2$};
  \node[ue,fill=violet!22,draw=violet!70!black] (v4) at (5.75,0.72) {UE$_4$};

  % tia toi
  \draw[gray!55,->] (1.06,1.85) -- (3.05,2.72);
  \draw[gray!55,->] (1.06,1.85) -- (3.05,1.05);
  % phan xa
  \draw[orange!80!black,->] (3.05,2.72) -- (1.80,3.10);
  \draw[blue!70!black,->]   (3.05,1.05) -- (1.98,0.75);
  % truyen qua
  \draw[teal!70!black,->,dashed]   (3.31,2.72) -- (5.28,3.02);
  \draw[violet!70!black,->,dashed] (3.31,1.05) -- (5.48,0.80);

  \node[font=\tiny,anchor=east,gray!45!black] at (3.05,3.74) {phía phản xạ};
  \node[font=\tiny,anchor=west,green!35!black] at (3.32,3.74) {phía truyền qua};

  \node[font=\scriptsize,anchor=north west] at (0,-0.30)
       {$\phi^T_n=\sqrt{\beta^T_n}\,e^{j\theta^T_n}$, \quad
        $\phi^R_n=\sqrt{\beta^R_n}\,e^{j\theta^R_n}$, \quad
        $\beta^T_n+\beta^R_n=1$};
\end{scope}

% =====================================================================
% Chi tiet mot phan tu (chung cho ca hai panel)
% =====================================================================
\begin{scope}[shift={(4.3,-3.60)}]
  \node[font=\scriptsize,anchor=south] at (2.9,1.62) {Chi tiết một phần tử STAR--RIS ở chế độ chia năng lượng (ES)};
  % tia toi
  \draw[gray!55,->,line width=0.7pt] (0.55,0.55) -- (2.55,0.55);
  \node[font=\tiny,anchor=north,gray!45!black] at (1.5,0.48) {tín hiệu tới};
  % phan tu
  \fill[green!20] (2.60,0.15) rectangle ++(0.30,0.80);
  \draw[green!55!black,line width=0.7pt] (2.60,0.15) rectangle ++(0.30,0.80);
  % phan xa
  \draw[orange!80!black,->,line width=0.7pt] (2.58,0.62) -- (1.58,1.12);
  \node[font=\tiny,anchor=south east,orange!70!black] at (1.66,1.08) {$\beta^R_n$ phản xạ};
  % truyen qua
  \draw[teal!70!black,->,line width=0.7pt,dashed] (2.92,0.55) -- (4.10,0.55);
  \node[font=\tiny,anchor=south,teal!60!black] at (3.72,0.62) {$\beta^T_n$ truyền qua};
  % thanh chia nang luong
  \draw[gray!60] (4.55,0.30) rectangle ++(1.70,0.45);
  \fill[teal!30] (4.55,0.30) rectangle ++(0.75,0.45);
  \fill[orange!30] (5.30,0.30) rectangle ++(0.95,0.45);
  \node[font=\tiny] at (4.93,0.525) {$\beta^T_n$};
  \node[font=\tiny] at (5.78,0.525) {$\beta^R_n$};
  \node[font=\tiny,anchor=north] at (5.40,0.24) {tổng bằng 1};
\end{scope}

\end{tikzpicture}
}
\caption{So sánh vùng phục vụ của RIS phản xạ truyền thống và STAR--RIS}
\label{fig:ch1-ris-star}
\figuresource{Tác giả xây dựng dựa trên \cite{Xu2021STARRIS, Mu2021STARAided} (2026)}
\end{figure}

\section{Học tăng cường sâu cho điều khiển liên tục}
\label{sec:drl}

\subsection{Từ học tăng cường đến DRL}
\label{ssec:tu-rl-den-drl}

Học tăng cường mô tả quá trình một tác tử tương tác với môi trường, quan sát trạng thái,
chọn hành động và nhận phần thưởng. Chuỗi tương tác thường được mô hình hóa bằng quá
trình quyết định Markov (Markov Decision Process, MDP). Trong nhiều bài toán điều khiển
cổ điển, trạng thái và hành động có thể rời rạc hóa để lưu giá trị trong bảng. Cách này
không phù hợp với bài toán STAR--RIS vì trạng thái chứa hàng trăm hệ số kênh thực và ảo,
còn hành động chứa hàng chục đến hàng trăm biến liên tục.

DRL thay bảng giá trị bằng mạng nơ-ron. Mạng có thể xấp xỉ hàm chính sách hoặc hàm giá
trị trên không gian liên tục có số chiều lớn. Tuy nhiên, DRL không loại bỏ nhu cầu thiết
kế bài toán. Nếu trạng thái thiếu thông tin hoặc hành động khó biểu diễn, mạng có thể
phải dành phần lớn năng lực để học các ràng buộc cơ bản thay vì học cách tối ưu hiệu
năng. Đây là lý do luận văn tách rõ ``mạng học'' và ``bộ giải mã vật lý''.

\subsection{Cấu trúc Actor--Critic}
\label{ssec:actor-critic}

Các thuật toán được sử dụng trong luận văn đều có thể nhìn dưới góc độ Actor--Critic.
Actor tạo hành động hoặc phân phối hành động từ trạng thái. Critic đánh giá chất lượng
của trạng thái hoặc cặp trạng thái--hành động. Hai mạng hỗ trợ nhau: Critic cung cấp tín
hiệu để Actor điều chỉnh, còn Actor tạo các hành động mà Critic cần đánh giá.

DDPG và TD3 là các thuật toán off-policy, nghĩa là có thể học từ dữ liệu được thu thập
bởi một chính sách hành vi không hoàn toàn giống chính sách hiện tại. Dữ liệu được lưu
trong bộ nhớ phát lại (Replay Buffer) và có thể được dùng lại nhiều lần, giúp tăng hiệu
quả sử dụng mẫu. PPO là thuật toán on-policy; dữ liệu mới được thu thập theo chính sách
hiện hành và được dùng trong một số vòng cập nhật trước khi bị thay bởi lô mới. Sự khác
nhau này ảnh hưởng đến độ ổn định, hiệu quả mẫu và cách thiết kế quy trình huấn luyện.

\subsection{DDPG}
\label{ssec:ddpg}

DDPG mở rộng ý tưởng gradient chính sách xác định sang không gian hành động liên tục
\cite{Lillicrap2015DDPG}. Actor sinh trực tiếp một véc-tơ hành động xác định. Critic ước
lượng giá trị $Q(s,a)$, tức mức phần thưởng tích lũy dự kiến khi thực hiện hành động $a$
tại trạng thái $s$. DDPG dùng mạng đích và bộ nhớ phát lại để ổn định học. Điểm yếu điển
hình là Critic có thể đánh giá quá cao một số hành động, khiến Actor bị đẩy về hướng
không thực sự tốt.

\subsection{TD3}
\label{ssec:td3}

TD3 được đề xuất để giảm sai số xấp xỉ của Actor--Critic bằng ba thay đổi chính
\cite{Fujimoto2018TD3}. Thứ nhất, hai Critic độc lập được dùng và mục tiêu lấy giá trị
nhỏ hơn để giảm hiện tượng đánh giá quá cao. Thứ hai, Actor cập nhật chậm hơn Critic để
Critic có thời gian thích nghi. Thứ ba, một nhiễu nhỏ được thêm vào hành động mục tiêu để
làm trơn giá trị và giảm sự phụ thuộc vào các đỉnh hẹp trong hàm Critic. Trong luận văn,
TD3 là một đại diện của nhóm off-policy chứ không được giả định là luôn tốt hơn DDPG hoặc
PPO.

\subsection{PPO}
\label{ssec:ppo}

PPO tối ưu chính sách bằng cách giới hạn mức thay đổi xác suất giữa chính sách mới và cũ
\cite{Schulman2017PPO}. Thay vì cho phép một bước gradient làm chính sách thay đổi quá
xa, PPO cắt tỷ số xác suất trong một khoảng lân cận. Cơ chế này đơn giản nhưng thường ổn
định trong thực nghiệm. Trong bài toán hành động liên tục, Actor của PPO tạo tham số của
phân phối Gaussian, sau đó lấy mẫu và đưa qua hàm Tanh để thu hành động nằm trong
$[-1,1]$. Một nghiên cứu gần đây đã áp dụng PPO cho bài toán tối đa hóa tổng tốc độ trong
mạng STAR--RIS hỗ trợ RSMA, cho thấy tính phù hợp của PPO với loại bài toán này
\cite{Meng2024PPOSTARRSMA}.

\begin{figure}[htbp]
\centering
\resizebox{0.92\textwidth}{!}{\begin{tikzpicture}[font=\scriptsize,>=Stealth,line width=0.5pt]
\node[box,text width=2.6cm,minimum height=1.0cm] (env) at (0,0)
     {Môi trường\\ STAR--RIS--RSMA};
\node[box,text width=2.4cm,minimum height=1.0cm] (st)  at (3.9,0)
     {Trạng thái\\ CSI chuẩn hóa};
\node[box,text width=2.4cm,minimum height=1.0cm] (ac)  at (7.6,0)
     {Actor\\ (chính sách)};
\node[box,fill=cChung!10,text width=2.6cm,minimum height=1.0cm] (dec) at (7.6,-2.2)
     {Bộ giải mã\\ hành động vật lý};
\node[box,text width=2.6cm,minimum height=1.0cm] (rw)  at (1.9,-2.2)
     {Tốc độ, QoS\\ và phần thưởng};

\draw[->] (env) -- (st);
\draw[->] (st)  -- (ac);
\draw[->] (ac)  -- (dec);
\draw[->] (dec) -- (rw);
\draw[->] (rw.north) -- (1.9,-1.05) -- (0,-1.05) -- (env.south);
\draw[->] (rw.east) -- (4.75,-2.2) -- (4.75,-1.15) -- (7.6,-1.15) -- (ac.south)
      node[pos=0.62,above,font=\tiny]{tín hiệu học};
\end{tikzpicture}
}
\caption{Vòng lặp DRL ở mức khái niệm trong bài toán phân bổ tài nguyên}
\label{fig:ch1-vong-lap-drl}
\figuresource{Tác giả xây dựng dựa trên \cite{Lillicrap2015DDPG, Fujimoto2018TD3, Schulman2017PPO} (2026)}
\end{figure}

\section{Các phương pháp tối ưu truyền thống}
\label{sec:toi-uu-truyen-thong}

\subsection{Tối ưu luân phiên}
\label{ssec:ao}

AO chia một bài toán nhiều khối biến thành các bài toán con. Tại một bước, một khối được
tối ưu trong khi các khối khác tạm thời giữ cố định. Sau đó chuyển sang khối tiếp theo và
lặp lại. Đối với STAR--RIS--RSMA, cách chia tự nhiên là khối công suất, khối phân bổ tốc
độ chung và khối STAR--RIS. Ưu điểm của AO là có thể tận dụng cấu trúc riêng của từng
khối. Nhược điểm là mỗi CSI cần nhiều vòng lặp và kết quả phụ thuộc nghiệm khởi tạo, thứ
tự cập nhật và các phép xấp xỉ cục bộ.

\subsection{Xấp xỉ lồi liên tiếp và cập nhật lân cận}
\label{ssec:sca}

Xấp xỉ lồi liên tiếp (Successive Convex Approximation, SCA) thay một bài toán khó bằng
chuỗi bài toán xấp xỉ dễ hơn quanh nghiệm hiện tại. Trong mã nguồn của luận văn, tên
AO--SCA được dùng cho một bộ giải cục bộ trong đó hai khối đơn hình được cập nhật bằng
chuyển khối khả thi và khối STAR--RIS dùng gradient sai phân hữu hạn kết hợp bước lân
cận. Vì vậy, luận văn mô tả đúng cơ chế triển khai thay vì đồng nhất bộ giải này với mọi
biến thể SCA trong tài liệu \cite{Scutari2017SCA}.

\subsection{Tìm kiếm theo lưới và căn pha giải tích}
\label{ssec:ao-grid}

AO--Grid thay gradient bằng một tập ứng viên rời rạc cho từng tọa độ. Với mỗi tọa độ, bộ
giải thử nhiều giá trị rồi giữ ứng viên cải thiện tốt nhất. Phương pháp này dễ kiểm soát
và có thể thực hiện các bước nhảy lớn, nhưng số phép đánh giá tăng nhanh theo số biến.
AnalyticalRIS là đối chiếu nhẹ hơn: sử dụng một quy tắc căn pha dựa trên kênh thay vì tối
ưu lặp đầy đủ. Nó có độ trễ thấp nhưng không tối ưu đồng thời tất cả biến RSMA và
STAR--RIS.

\section{Các nghiên cứu liên quan}
\label{sec:nghien-cuu-lien-quan}

Các công trình liên quan có thể chia thành bốn nhóm. Nhóm thứ nhất xây dựng nền tảng RSMA
và chỉ ra khả năng quản lý nhiễu mềm giữa SDMA và NOMA \cite{Mao2018RSMABridging,
Mao2022RSMAFundamentals, Joudeh2016RSMASumRate}. Nhóm thứ hai nghiên cứu RIS và
STAR--RIS, từ mô hình phản xạ truyền thống đến mô hình đồng thời truyền và phản xạ cùng
các ràng buộc pha \cite{Xu2021STARRIS, Mu2021STARAided, Wu2019IRSBeamforming,
DiRenzo2020SmartRadio, Wu2021IRSTutorial, Liu2021STARCoupledPhase,
Wang2022STAROptimization}. Nhóm thứ ba kết hợp RSMA với RIS hoặc STAR--RIS để tối ưu tốc
độ, bảo mật hoặc hiệu năng hệ thống \cite{Li2022RSMARISInterplay,
Aboumahmoud2024RISRSMASurvey, Ge2024RSMAMassiveMIMO, Liu2024DiscretePhaseRSMARate,
Hashempour2024SecureSTARRS, Maghrebi2024CooperativeActiveSTAR}. Nhóm thứ tư dùng học tăng
cường cho STAR--RIS hoặc RSMA, trong đó có PPO và các phương pháp học lai
\cite{Meng2024PPOSTARRSMA, Zhong2022HybridRLSTAR}.

Bảng~\ref{tab:ch1-nghien-cuu} tổng hợp một số công trình tiêu biểu theo mục tiêu và hạn
chế có liên quan trực tiếp đến luận văn. Bảng được tác giả tự tổng hợp từ nội dung công
bố, không sao chép bảng từ tài liệu khác.

\begin{table}[htbp]
\centering
\caption{Tổng hợp một số hướng nghiên cứu liên quan}
\label{tab:ch1-nghien-cuu}
\footnotesize
\begin{tabularx}{\textwidth}{@{}l l Y@{}}
\toprule
\textbf{Công trình} & \textbf{Hệ thống và phương pháp} & \textbf{Điểm liên quan và giới hạn so với luận văn} \\
\midrule
Mao và cộng sự \cite{Mao2018RSMABridging, Mao2022RSMAFundamentals}
& RSMA nhiều người dùng; phân tích và tối ưu
& Cung cấp nền tảng luồng chung và luồng riêng; không tập trung STAR--RIS và biểu diễn hành động DRL. \\
\addlinespace
Xu, Mu và cộng sự \cite{Xu2021STARRIS, Mu2021STARAided}
& STAR--RIS; mô hình hệ thống ES, MS, TS
& Cung cấp mô hình truyền và phản xạ; không giải bài toán DRL của luận văn. \\
\addlinespace
Li và cộng sự \cite{Li2022RSMARISInterplay}
& RIS kết hợp RSMA; phân tích tương tác
& Làm rõ tiềm năng kết hợp hai công nghệ; không tập trung cấu trúc SISO và bộ giải mã học. \\
\addlinespace
Meng và cộng sự \cite{Meng2024PPOSTARRSMA}
& STAR--RIS hỗ trợ RSMA; PPO
& Chứng minh PPO có thể dùng cho bài toán liên tục; luận văn mở rộng đánh giá sang DDPG, TD3 và nhóm AO trên cùng mô hình. \\
\addlinespace
Zhong và cộng sự \cite{Zhong2022HybridRLSTAR}
& STAR--RIS; học tăng cường lai
& Khai thác cấu trúc điều khiển STAR--RIS; mục tiêu và mô hình khác RSMA SISO của luận văn. \\
\addlinespace
Ge và cộng sự \cite{Ge2024RSMAMassiveMIMO}
& STAR--RIS, massive MIMO, RSMA; tối ưu chung
& Hệ nhiều anten và mục tiêu khác; không phản ánh đặc điểm không có beamforming của SISO. \\
\bottomrule
\end{tabularx}
\tablesource{Tác giả tổng hợp từ các công trình trích dẫn (2026)}
\end{table}

\section{Khoảng trống nghiên cứu và định vị của luận văn}
\label{sec:khoang-trong}

Tổng quan cho thấy có nhiều nghiên cứu về RSMA, STAR--RIS và DRL, nhưng bài toán của luận
văn có một giao điểm cụ thể: hệ SISO không có phân tách không gian, trong khi hành động
vẫn chứa nhiều công suất luồng riêng và hàng loạt biến STAR--RIS. Nếu mạng DRL được yêu
cầu học toàn bộ miền khả thi từ đầu, nó phải đồng thời học rằng nhiều luồng riêng có thể
gây nhiễu mạnh, học cách chọn UE ưu tiên, học cách chia luồng chung, học năng lượng hai
phía và học pha cho hàng chục hoặc hàng trăm phần tử.

Khoảng trống mà luận văn tập trung không phải tuyên bố rằng chưa từng có DRL cho
STAR--RIS--RSMA. Thay vào đó, nghiên cứu đặt trọng tâm vào cách biểu diễn hành động: liệu
có thể dùng phân tích vật lý của SISO--RSMA và căn pha để tạo một miền hành động thuận
lợi hơn, sau đó dùng nhiều thuật toán DRL khác nhau để học trên miền này hay không. Cách
định vị này phù hợp với kết quả thực nghiệm khi TD3, DDPG và PPO không có một thuật toán
luôn thắng tuyệt đối, nhưng cả nhóm DRL có thể tạo quyết định rất nhanh so với AO.

Luận văn vì vậy không xem một thuật toán học là đóng góp duy nhất. Đóng góp nằm ở sự kết
hợp giữa mô hình vật lý, phân tích cấu trúc, bộ giải mã hành động và giao thức đánh giá.
Các thuật toán TD3, DDPG và PPO được dùng như ba cơ chế học khác nhau trên cùng nền tảng
để kiểm tra mức độ tổng quát của cách biểu diễn.

\section{Kết luận chương}
\label{sec:ch1-ket-luan}

Chương~\ref{chap:co-so-ly-thuyet} đã trình bày nền tảng về đa truy cập, RSMA, RIS và
STAR--RIS, DRL cùng các phương pháp tối ưu lặp. Điểm quan trọng được rút ra là bài toán
SISO STAR--RIS--RSMA không chỉ khó vì số chiều lớn mà còn có cấu trúc nhiễu đặc thù do
thiếu phân tách không gian. Chương tiếp theo xây dựng mô hình toán học đúng với mã nguồn
và phân tích chi tiết cấu trúc này để hình thành cơ sở cho bộ giải mã hành động ở
Chương~\ref{chap:phuong-phap}.
